\chapter{哈尔测度}

在本章和下一章中，当我们说到函数（分别地，测度）时，应理解为实函数或复函数（分别地，测度）；若 T 是局部紧空间，记号 \(\mathcal{K}(T)\) 将表示空间 \(\mathcal{K}_{\mathbf{R}}(T)\) 或空间 \(\mathcal{K}_{\mathbf{C}}(T)\)；对于记号 \(\overline{\mathcal{K}(T)}\)、\(\mathcal{C}(T)\)、\(L^{p}(T,\mu)\)、\(\mathcal{M}(T)\) 等也是如此。当然，在涉及多个函数、测度或向量空间的情形，所得结果在这些函数、测度或向量空间全为实的或全为复的时均成立。总是假定空间 \(\overline{\mathcal{K}(T)}\) 配备一致收敛拓扑，空间 \(\mathcal{C}(T)\) 配备紧收敛拓扑，空间 \(\mathcal{K}(T)\) 配备直接极限拓扑，其定义在第六章开头回顾。记号 \(\mathcal{K}_{+}(T)\) 表示满足 \(\geqslant0\) 的 \(\mathcal{K}(T)\) 函数集合。若 \(A\subset T\)，则 \(\varphi_{A}\) 总表示 A 的示性函数。若 \(t\in T\)，则 \(\varepsilon_{t}\) 表示在点 t 处质量为 +1 的正测度。

所有局部凸空间均假定为豪斯多夫空间。

除非明确另有说明，我们将 e 记为所考虑的所有群的单位元。

